% =====================================================================================================
\section{Setting and notation}\label{sec:setting}

This section fixes the data, the model, the prior and the posterior, and names the few objects of the method that the
theorems talk about. The method is described in full in the methodology section; here we recall only what the statements
need. Every symbol is listed in Appendix~\ref{app:notation}.

\subsection{Notation}\label{sec:notation}

For an integer $k\ge1$ we write $\set{k}:=\{1,\dots,k\}$. Vectors are columns. The Euclidean norm of a vector $v$ is
$\norm{v}$ and its $\ell_1$ norm is $\norm{v}_1$. An identity matrix is written $I_n$, or $I$ when its size is clear,
and $\one_n$ (or $\one$) is the vector of ones. The smallest eigenvalue of a symmetric matrix $B$ is $\lmin(B)$. The
number of elements of a set $S$ is $\abs{S}$. Logarithms are natural ($\log$) or binary ($\log_2$);
$\lceil\cdot\rceil$ and $\lfloor\cdot\rfloor$ round up and down; $a\vee b:=\max(a,b)$; and $\ind{E}$ is the indicator of
an event or condition $E$. A star marks a true parameter value; the superscript $\mathrm r$ marks raw, unstandardised
data.

\subsection{Data and model}\label{sec:model}

We observe a response vector $\yr\in\R^n$ and $\praw$ predictors. The method first removes every predictor that is
constant in the sample. A constant predictor is a multiple of $\one_n$, so its effect is absorbed by the intercept. We
therefore write the model with the remaining predictors only: $\Xr$ is the $n\times p$ matrix of the $p$ \emph{retained}
predictors, indexed by $\set{p}$. Everything the method computes, including its prior, uses these $p$ predictors only.
The raw count $\praw$ appears only in Section~\ref{sec:distribution} and in the computing conventions.

The data follow a Gaussian linear model with an intercept:
\begin{equation}\label{eq:model}
  \yr=\alpha^*\one_n+\Xr\bstar+\varepsilon,\qquad \varepsilon\sim\Normal_n(0,\sigs^2I_n),
\end{equation}
with $\sigs^2>0$ and $\varepsilon$ independent of $\Xr$. In words: each response is a linear function of the predictors
plus independent Gaussian noise of variance $\sigs^2$.

A \emph{model} is a set $S\subseteq\set{p}$ of predictors. The \emph{true model} is the set of predictors with a nonzero
coefficient, $\Sz:=\supp\bstar=\{j:\beta^*_j\neq0\}$, and $\sz:=\abs{\Sz}$ is its size. The members of $\Sz$ are the
\emph{true predictors}; the others are \emph{null predictors}.

We consider two kinds of design.
\begin{itemize}[leftmargin=1.5em]
  \item \emph{Fixed design.} The matrix $\Xr$ is deterministic, and probability refers to the noise $\varepsilon$ only.
  \item \emph{Random Gaussian design.} The rows of $\Xr$ are independent and identically distributed (i.i.d.)
    $\Normal_p(m_n,\Sig)$, and probability refers to the joint law of $(\Xr,\varepsilon)$. The subscript $n$ recalls
    that this law may change with $n$. When the rows are random, we write $W=(W_1,\dots,W_p)\sim\Normal_p(0,\Sig)$ for a
    centred copy of one row; for $S\subseteq\set{p}$, $W_S$ is the sub-vector indexed by $S$ and $\Sigma_{n,S}$ the
    principal sub-matrix of $\Sig$ indexed by $S$.
\end{itemize}
No theorem assumes a particular correlation structure, such as equicorrelation; such structures appear only as
examples.

\subsection{What the method sees: standardised data}\label{sec:standardise}

The method centres and scales every predictor and the response. Let $P^\perp_{\one}:=I_n-n^{-1}\one\one^\top$ be the
centring matrix and $D$ the $p\times p$ diagonal matrix of the sample standard deviations (divisor $n-1$) of the columns
of $\Xr$. The method works with
\begin{equation}\label{eq:standardise}
  X:=P^\perp_{\one}\Xr D^{-1},\qquad
  y:=P^\perp_{\one}\yr/s_y,\qquad s_y^2:=\norm{P^\perp_{\one}\yr}^2/(n-1).
\end{equation}
Each column $x_j$ of $X$ has mean zero and $\norm{x_j}^2=n-1$, and $y$ has mean zero and $\norm{y}^2=n-1$.

In terms of the standardised predictors, the centred response is
\[
  P^\perp_{\one}\yr=X\tstar+P^\perp_{\one}\varepsilon,\qquad \tstar:=D\bstar .
\]
We call $\tstar$ the \emph{standardised coefficients}; they are in the units of the raw response. The \emph{mean signal}
is $\mu:=X\tstar\in\R^n$, and $\tmin:=\min_{j\in\Sz}\abs{\theta^*_j}$ is the smallest true standardised coefficient.
Standardising does not move the support: $\supp\tstar=\Sz$. For a fixed design $\tstar$ is deterministic; for a random
design it is random through $D$, and a theorem that places its conditions on $\tstar$ says so.

For a model $S$: $X_S$ is the $n\times\abs{S}$ matrix of the columns in $S$; $G_S:=X_S^\top X_S$ is their Gram matrix;
$\Rhat_S:=G_S/(n-1)$ is their sample correlation matrix; and $P_S$ is the orthogonal projection onto the span of the
columns of $X_S$, with $P_\emptyset:=0$.

\subsection{The prior, as coded}\label{sec:prior}

The method fits a working Gaussian linear model to the standardised data and puts a prior on its three unknowns: which
predictors are in the model, their coefficients, and the noise variance. Nothing in the prior is tuned by the user.

\paragraph{Working model.} Given a model $S$, a coefficient vector $\beta_S\in\R^{\abs{S}}$ and a variance
$\sigma^2>0$, the working likelihood is
\begin{equation}\label{eq:working}
  y\mid S,\beta_S,\sigma^2\ \sim\ \Normal_n(X_S\beta_S,\sigma^2I_n).
\end{equation}
Here $\beta_S$ and $\sigma^2$ are working parameters, not the true values. The intercept is removed by centring and has no
prior of its own. The theorems are proved under the true model~\eqref{eq:model}, not under this working model.

\paragraph{Model space.} The method considers every model with at most $K$ predictors:
\begin{equation}\label{eq:modelspace}
  \Om:=\{S\subseteq\set{p}:\abs{S}\le K\},\qquad K:=\min\big(p,\lfloor n/2\rfloor\big).
\end{equation}

\paragraph{Model prior.} The prior probability of a model depends only on its size:
\begin{equation}\label{eq:modelprior}
  \text{prior probability of }S\;=\;\frac{w(\abs{S})}{W_K},\qquad w(k):=\frac{1}{2\binom{2p-1}{k}},\qquad W_K:=\sum_{k=0}^{K}\binom{p}{k}w(k),\qquad
  S\in\Om .
\end{equation}
This is the beta-binomial prior $\BB(1,p)$ \citep{ScottBerger2010}, restricted to models of size at most $K$: a
probability $q$ is drawn from $\BetaDist(1,p)$, and, given $q$, each predictor enters the model independently with
probability $q$. Its prior mean model size is about one. Adding one predictor to a model of size $k$ multiplies the
prior weight by $(k+1)/(2p-1-k)$; that is, one more predictor costs about $\log\{2p/(k+1)\}$ in log prior odds. The
methodology section works through the case $p=3$.

\paragraph{Coefficients and noise variance.} Given $S$ and $\sigma^2$, the coefficients have the ridge-stabilised
$g$-prior \citep{Zellner1986} with $g=n$, and the noise variance has an inverse-gamma prior with shape $a_0$ and scale
$b_0$:
\begin{equation}\label{eq:slab}
  \beta_S\mid S,\sigma^2\ \sim\ \Normal\big(0,\;n\sigma^2(G_S+I)^{-1}\big),\qquad
  \sigma^2\ \sim\ \IG(a_0,b_0),\quad a_0=b_0=0.01.
\end{equation}

\subsection{The posterior}\label{sec:posterior}

Integrating $\beta_S$ and $\sigma^2$ out of~\eqref{eq:working}--\eqref{eq:slab} gives the marginal likelihood of each
model in closed form:
\begin{equation}\label{eq:marglik}
  m(S)=C_n\Big(\frac{\det H_S}{\det(G_S+I)}\Big)^{-1/2}\,(2b_0+\SSq_S)^{-\nu},\qquad
  H_S:=(n+1)G_S+I,\quad \SSq_S:=y^\top y-n\,z_S^\top H_S^{-1}z_S,
\end{equation}
where $z_S:=X_S^\top y$, $\nu:=a_0+n/2$ and $C_n:=(2b_0)^{a_0}\Gamma(a_0+n/2)/\{\pi^{n/2}\Gamma(a_0)\}$ (here $\pi$ is the
number $3.14\ldots$). In words, $\SSq_S$ is a shrunken residual sum of squares of model $S$, and the determinant ratio is
a penalty for model size (an Occam factor): in a well-conditioned model, each extra predictor multiplies $m(S)$ by about
$(n+1)^{-1/2}$.

The \emph{posterior over models} is
\begin{equation}\label{eq:posterior}
  \pis{S}=\frac{w(\abs{S})\,m(S)}{W_K\,Z},\qquad Z:=\frac{1}{W_K}\sum_{V\in\Om}w(\abs{V})\,m(V),\qquad S\in\Om .
\end{equation}
The normalising constant $Z$ is the \emph{evidence}: the marginal likelihood of the standardised data under the working
model and the prior. The \emph{log score} of a model is $F(S):=\log\{w(\abs{S})/W_K\}+\log m(S)=\log\pis{S}+\log Z$: the log
posterior, up to the additive constant $\log Z$. The method's target is exactly this posterior. Because $y$ is standardised, the posterior over models
does not change if $\yr$ is shifted or rescaled.

\paragraph{Posterior summaries.} These are properties of the posterior, not of any algorithm:
\begin{itemize}[leftmargin=1.5em]
  \item the \emph{posterior inclusion probability} of predictor $j$, $\PIP_j:=\sum_{S\ni j}\pis{S}$;
  \item the \emph{highest posterior model} $\HPM:=\argmax_{S\in\Om}\pis{S}$;
  \item the \emph{median probability model} $\MPM:=\{j:\PIP_j>1/2\}$ \citep{BarbieriBerger2004};
  \item the \emph{posterior miss} $1-\pis{\Sz}$, the posterior mass of the wrong models (with $\pis{\Sz}:=0$ if
    $\Sz\notin\Om$).
\end{itemize}
The next lemma is used throughout: a bound on the posterior miss controls all three summaries at once.

\begin{lemma}\label{lem:summaries}
For every data set, $\max_{j\in\set{p}}\abs{\PIP_j-\ind{j\in\Sz}}\le1-\pis{\Sz}$. If $\pis{\Sz}>1/2$, then
$\HPM=\MPM=\Sz$.
\end{lemma}

\begin{proof}
For $j\in\Sz$, $1-\PIP_j$ is the mass of models that miss $j$, all of which differ from $\Sz$; for $j\notin\Sz$,
$\PIP_j$ is the mass of models that contain $j$, all of which differ from $\Sz$. If $\pis{\Sz}>1/2$, then $\Sz$ beats
every other model, and the first part gives $\PIP_j>1/2$ exactly for $j\in\Sz$.
\end{proof}

So every inclusion probability is within $1-\pis{\Sz}$ of its ideal value: $1$ for a true predictor and $0$ for a null
one.

\subsection{A run of the method, and what it returns}\label{sec:method}

The method is a sequential Monte Carlo sampler over $\Om$, combined with a local search. Four of its objects appear in
the theorems: the start law, the local search, the reported model, and the estimated summaries.

\paragraph{A run.} A \emph{run} is one call of the software on a data set, with a user-supplied random \emph{seed}. When
$\Om$ is small enough, the method enumerates it and computes the posterior exactly. Otherwise it runs the sampler. Before
the final sampler run it performs an internal \emph{calibration}: a few short pilot runs that fix the sequence of
intermediate distributions and the effort level, so that these settings are frozen before the final run's random numbers
are drawn. A run either \emph{produces output}, ending with status ``exact'' (enumeration) or ``completed'' (sampler), or
it stops without estimates, for example when its computing allocation is exhausted. We also say that a run that produces
output \emph{completes}.

\paragraph{The start law.} The method partitions the predictors into blocks $C_1,\dots,C_J$ of at most $16$ predictors
each, by clustering on sample correlations; the theory uses only that the partition is a function of $X$ with blocks of
size at most $16$. Within each block the method scores every subset exactly. The \emph{start law} is the posterior with
all interactions between blocks ignored:
\begin{equation}\label{eq:startlaw}
  \pi_0(S)\ \propto\ w(\abs{S})\prod_{b=1}^{J}\frac{m(S\cap C_b)}{m(\emptyset)},\qquad S\in\Om.
\end{equation}
It can be computed and sampled exactly, and the sampler starts from it. The \emph{start mode} $\tSz$ is a maximiser of
$\pi_0$ over $\Om$.

\paragraph{Local search.} Two models are \emph{neighbours} if one comes from the other by deleting one predictor, adding
one (while the size is below $K$), or swapping one member for one non-member. A local search scans all neighbours of the
current model and moves to the best one if its log score is strictly higher. It stops at the first scan with no strict
improvement, or when it reaches its move cap, whichever comes first. After the sampler, the method runs two local
searches: one from the best model found so far, and one from the start mode, with move cap $\ell_n+\abs{\tSz}$, where
$\ell_n:=\lceil\log_2(n+p+1)\rceil$.

\paragraph{The reported model and certification.} The method records every model it scores, in the sampler and in the
local searches. The \emph{reported model} $\Shat$ is the best-scoring recorded model. The method \emph{certifies}
$\Shat$ (it sets its local-optimality flag) in two cases: (a) it enumerated $\Om$, so that $\Shat$ maximises $F$ over
$\Om$; or (b) $\Shat$ ends a local search whose last scan found no strict improvement, so that no neighbour of $\Shat$
has a higher log score.

\paragraph{Estimated summaries.} The sampler runs four independent copies, called \emph{islands}. Island $r$ produces
estimated inclusion probabilities $\PIPhat_{r,j}$ and an estimate $\Zhat_r$ of the evidence $Z$. The method pools the
islands with equal weights,
\[
  \PIPhat_j:=\frac14\sum_{r=1}^4\PIPhat_{r,j},\qquad \Zhat:=\frac14\sum_{r=1}^4\Zhat_r,
\]
and reports the estimated median probability model $\MPMhat:=\{j:\PIPhat_j>1/2\}$.

\subsection{Asymptotic conventions}\label{sec:asymptotics}

All results are asymptotic in $n\to\infty$, along a sequence of problems: $p=p_n$, the design, $\bstar$ and $\Sz$ may
change with $n$. A quantity is \emph{fixed} if it does not depend on $n$; every constant introduced in an assumption is
fixed. Generic positive constants that depend only on fixed quantities are written $C$ and $c$ in prose; a constant that
appears in a displayed result has its own name.

For random variables $V_n$ and positive numbers $a_n$: $V_n=\OP(a_n)$ means that for every $\epsilon>0$ there is a
constant $C$ with $\P(\abs{V_n}>Ca_n)\le\epsilon$ for all large $n$; $V_n=\oP(a_n)$ means $V_n/a_n\to0$ in probability.
\emph{With probability tending to one} means that the probability of the event tends to one. For positive sequences,
$a_n=O(b_n)$ (also $a_n\lesssim b_n$) means $a_n\le Cb_n$ for a constant $C$; $a_n=o(b_n)$ means $a_n/b_n\to0$;
$a_n\asymp b_n$ means $a_n\lesssim b_n$ and $b_n\lesssim a_n$; and $a_n\gg b_n$ means $b_n/a_n\to0$.

\subsection{Standing assumptions and conventions}\label{sec:standing}

Most theorems use the following two assumptions; each theorem lists the ones it uses. Three conventions about the
computation are used only where a theorem is about what the software does.

\begin{assumption}{GL}{Gaussian linear model}
The data follow model~\eqref{eq:model}, with a fixed noise variance $\sigs^2>0$ and noise independent of the design.
\rulesout{heavy-tailed or heteroscedastic noise, and noise that depends on the predictors.}
\end{assumption}

\begin{assumption}{DIM}{Dimension}
$2\le p\le n^A$ for a fixed $A\ge1$.
\rulesout{a number of predictors growing faster than every power of $n$, such as $p=e^{n^a}$. The case $p\gg n$ is
allowed and is the main case.}
\end{assumption}

The conventions are not statistical assumptions: they fix what ``the method'' means and idealise the computer.

\begin{convention}{C1}{The method}
Statements about what the method returns assume exact arithmetic (scores and comparisons computed without rounding
error), the default configuration (in particular, the local search from the start mode is switched on), and valid inputs
and seed. Unless a statement says otherwise, it holds for every valid seed and every outcome of the internal
calibration.
\end{convention}

\begin{convention}{C2}{Delivery}
Used only by the clauses ``the run always produces output''. The computer does exact arithmetic, never overflows and
never fails. The method runs in its default configuration, with ideal random numbers. The user grants at least the
memory and time that the method's own requirement formulas, polynomial in $(n,\praw)$, ask for. Convention
\as{C2p} is the same with the code's two fixed integer limits on the local-search move caps lifted. Appendix~\ref{app:conventions}
gives the full lists.
\end{convention}

\begin{convention}{C3}{Sampler}
Used only by the results on the estimated summaries. Exact arithmetic; ideal independent uniform random numbers, with the
final run's random numbers independent of the data and of the calibration; exact draws from the start law; and Markov
moves that do exactly what their definitions say. Appendix~\ref{app:conventions} gives the full list.
\end{convention}
